\documentclass[12pt, a4paper]{article}
\usepackage[utf8]{inputenc}
\usepackage{amsmath, amssymb, amsthm}
\usepackage{geometry}
\geometry{margin=1in}
\usepackage{array}
\usepackage{booktabs}

\begin{document}

\section{Solidity, Containment, and the Origin of Weight}

\subsection{What Solidity Is in GPT}

Solidity is not an intrinsic property of atoms. It is a \textbf{coherence state of the graviton field}. A solid object is a region where graviton flow is contained — closed-loop internal circulation, no net external flow. This containment \textit{is} solidity.

\subsection{The Absorption/Ejection Cycle in a Solid}

Gravitons are absorbed, processed, and re-emitted. In a solid, the outflow circulates internally rather than escaping. Self-repulsion creates a stable internal pressure profile.

\subsection{The Normal Force as Containment-Containment Boundary}

When two solids meet, their containment fields cannot interpenetrate without reorganization. Self-repulsion creates a boundary pressure. This is the normal force.

\subsection{Weight as Asymmetric External Pressure on a Contained Field}

Earth's graviton field flows inward, creating an anisotropic pressure gradient. A solid object's containment field resists this gradient. The net force is weight.

\subsection{Why Weight Scales with Mass}

Containment field strength scales with graviton processing rate, which scales with atomic count. Thus \( F_{\text{weight}} \propto m \).

\subsection{The Equivalence Principle from the Medium Perspective}

Both gravitational weight (anisotropic ambient density) and acceleration weight (motion-induced anisotropy) produce identical containment field boundary stress. The EP is not a coincidence — it is a consequence of weight being a containment field phenomenon.

\section{Orbital Mechanics as Spatiotemporal Resonance}

\subsection{The Intake/Expulsion Cycle in a Non-Uniform Object}

Real objects are non-uniform. Absorption and ejection rates vary across the surface.

\subsection{Spin Creates Temporal Modulation (The Object's Pulse)}

When the object rotates, the non-uniform surface is presented to the inflow in a time-dependent pattern. The object pulses at frequency \( f_{\text{obj}} = k\Omega \).

\subsection{The Graviton Field's Pressure Oscillation}

The central mass's absorption/ejection cycle creates field oscillations. Shell stratification produces frequency bands \( f_{\text{field}}(r_n) \).

\subsection{Resonance Condition for Stable Orbit}
\[
f_{\text{obj}} = m \cdot f_{\text{field}}(r_n) \quad \text{or} \quad m \cdot f_{\text{obj}} = f_{\text{field}}(r_n)
\]

Stable orbits exist only where this condition is satisfied.

\subsection{Why Orbital Radii Are Quantized}

The field's frequency spectrum is discrete (core pulsation + shell boundaries). Only discrete radii satisfy resonance.

\subsection{Orbital Velocity as Consequence, Not Cause}

Velocity is derived from pressure gradient balance at the resonant radius, not the independent variable.

\subsection{Spin-Orbit Coupling and Tidal Locking}

Because \( f_{\text{obj}} = k\Omega \), spin and orbit are linked. When \( \Omega = \omega_{\text{orb}} \), tidal locking occurs.

\subsection{Why Spin and Orbit Are Aligned (Prograde vs. Retrograde)}

Retrograde rotation produces a pulse pattern inverted relative to field oscillation — destructive interference, decoherence, orbital decay.

\subsection{Scaling Invariance: From Planetary Orbits to Atomic Shells}

The resonance mechanism scales down to the atom: nucleus as central oscillator, electron as spinning non-uniform structure, discrete radii as electron shells. Chemical bonding is shared resonance corridor between atoms.

\subsection{Answering the Opening Question: Why Structure Survives}

Uncertainty is not fundamental. It is undersampling of a deterministic medium. Structure survives because resonance is enforced by the field.

\section{Light Deflection as Refraction in a Pressure Field}

\subsection{The Refractive Index of the Graviton Medium}
\[
n(r) = 1 + \frac{2GM}{c^2 r}
\]

\subsection{Fermat's Principle in a Density Gradient}

Light follows \( \delta \int n(r) \, ds = 0 \). The deflection angle:
\[
\alpha = \frac{4GM}{c^2 b}
\]

For the Sun: \( \alpha = 1.75 \) arcseconds — same as GR.

\subsection{Why Light Bends Toward the Mass}

Light bends because the medium's density gradient creates a refractive index gradient. Same mechanism as any wave bending in a non-uniform medium.

\subsection{Refraction vs. Curvature: A Critical Distinction}

GR: geodesics in curved geometry. GPT: refraction in varying medium density. Same numbers. Different ontology. Only GPT provides a mechanism.

\section{Inertia and Momentum as Coherence Retention}

\subsection{Momentum as a Carved Passage (The Momentum Corridor)}

Moving object absorbs more gravitons on leading edge, creating a vacancy trail. The medium retains a structured flow pattern — a low-impedance channel. This structured flow \textit{is} momentum.

\subsection{Why Momentum Is Conserved}

The medium's structured flow is conserved. Noether's theorem describes this consequence, but the cause is the medium itself.

\subsection{Inertia as Resistance to Reorientation}

Changing velocity requires disrupting the existing corridor and carving a new one. The medium resists because the existing corridor is a lower-energy configuration.

\subsection{Why Inertia Scales with Mass}

Mass is coupling strength to the medium. Larger mass = stronger corridor = greater inertia.

\section{Time as Process: The N/Γ Cycle}

\subsection{Time as a Count of Cycles}
\[
\tau = \frac{N}{\Gamma}
\]

\subsection{Gravitational Time Dilation}

Higher ambient density increases load factor \( N \) faster than encounter rate \( \Gamma \):
\[
\tau(r) = \tau_0 \cdot \frac{1}{\sqrt{1 - 2GM/(c^2 r)}}
\]

\subsection{Velocity Time Dilation}
\[
\tau(v) = \tau_0 \cdot \gamma(v) = \tau_0 \cdot \frac{1}{\sqrt{1 - v^2/c^2}}
\]

\section{Charge as Stabilized Bias}

\subsection{The Exchange Kernel and Bias}

For a symmetric structure, \( \int K(\Omega) \Omega \, d\Omega = 0 \). A charged structure has non-zero first moment — directional bias.

\subsection{Why Two Signs Exist}
\[
\mathbf{B} = \pm B_0 \hat{\mathbf{s}}
\]
Two stable orientations of the bias vector.

\subsection{Why Opposite Charges Attract}

Outflow from positive feeds into inflow of negative. Shared corridor reduces medium stress. Ambient pressure pushes them together.

\subsection{Why Like Charges Repel}

Both outflow or both inflow conflict in the same region. Self-repulsion or vacancy conflict creates high pressure between them.

\section{The Electric Field as Bias-Potential Gradient}

\subsection{The Scalar Bias Potential}
\[
\nabla^2 \psi = -\alpha \sigma_q
\]

\subsection{The Electric Field as Gradient}
\[
\mathbf{e} = -\nabla \psi
\]

\subsection{Inverse-Square from Flux Dilution}
\[
e(r) = \frac{\alpha q}{4\pi r^2} \hat{\mathbf{r}}
\]

\subsection{Permittivity as Medium Compressibility}
\[
\epsilon_0 = \frac{1}{\alpha}
\]

\section{Magnetism as Curl Organization from Moving Bias}

\subsection{The Vector Circulation Potential}
\[
\nabla^2 \mathbf{A} = -\beta \mathbf{J}_q
\]

\subsection{The Magnetic Field as Curl}
\[
\mathbf{b} = \nabla \times \mathbf{A}
\]

\subsection{Permeability as Transverse Curl Stiffness}
\[
\mu_0 = \beta
\]

\section{Constants as Medium Properties}

\subsection{\( \epsilon_0 \) as Bias Compressibility}

How much bias potential gradient is generated per unit bias source.

\subsection{\( \mu_0 \) as Transverse Curl Stiffness}

How much the medium resists shear from moving bias.

\subsection{\( c^2 = 1/(\epsilon_0 \mu_0) \)}

Necessary relation for a medium supporting both longitudinal and transverse waves at same speed.

\section{The Lorentz Force as Medium Response}

\subsection{Electric Force}
\[
\mathbf{F}_e = q \mathbf{e}
\]

\subsection{Magnetic Force}

Lowest-order rotationally consistent form:
\[
\mathbf{F}_m = q \mathbf{v} \times \mathbf{b}
\]

\subsection{Complete Lorentz Force}
\[
\mathbf{F} = q(\mathbf{e} + \mathbf{v} \times \mathbf{b})
\]

\subsection{No Magnetic Monopoles}

Because \( \mathbf{b} = \nabla \times \mathbf{A} \), \( \nabla \cdot \mathbf{b} = 0 \) identically.

\section{Electromagnetic Waves as Oscillating Medium Disturbances}

\subsection{Coupled Wave Equations}
\[
\nabla^2 \psi - \frac{1}{c^2} \frac{\partial^2 \psi}{\partial t^2} = 0, \quad
\nabla^2 \mathbf{A} - \frac{1}{c^2} \frac{\partial^2 \mathbf{A}}{\partial t^2} = 0
\]

\subsection{What Is Oscillating?}

The medium's bias potential \( \psi \) and vector potential \( \mathbf{A} \). Light is a wave in the graviton medium.

\subsection{Why Light Has No Rest Frame}

The medium's rest frame is inaccessible to massive objects because their own motion is also mediated by the medium.

\subsection{Photons as Coherent Wave Packets}

Quantization emerges from the medium's coherence requirements. \( E = \hbar \omega \) where \( \hbar \) is the minimal action unit of the medium.

\section{Light Deflection as Refraction in a Pressure Field}

\subsection{Refractive Index from Density Perturbation}
\[
n(r) = 1 + \frac{2GM}{c^2 r}
\]

\subsection{Deflection Angle}
\[
\alpha = \frac{4GM}{c^2 b}
\]

\subsection{Shapiro Delay as Slowing in Denser Medium}
\[
\Delta t = \frac{2GM}{c^3} \ln \left( \frac{4r_1 r_2}{b^2} \right)
\]

\subsection{Gravitational Lensing as Refractive Focusing}

Multiple images, Einstein rings, arcs arise from refractive index gradients, not curved spacetime.

\subsection{Testable Prediction}

Small wavelength-dependent dispersion (\( \alpha(\lambda) \) varies) — not present in GR.

\end{document}